% GENERATED FILE. Edit canonical lessons, then run npm run generate:books.
% canonical-source-hash: fnv1a64:d022818b101216e2
% canonical-lessons: SA-C476-saktih, SA-C476-pramsuh, SA-C476-pustah, SA-C476-akrtih, SA-C476-posayati

\chapter{Strength and Build}
\label{ch:sa-strength-and-build}
\hlchaptermodality{476}

\begin{chapteropening}
\textbf{By the end of this chapter:} \emph{I can describe how strong or well-built someone is.}

Complete the last lesson of chapter 476: I can describe how strong or well-built someone is.
\end{chapteropening}

\section[\texorpdfstring{\sk{शक्तिः} (śaktiḥ)}{शक्तिः (śaktiḥ)}]{\sk{शक्तिः} (śaktiḥ) — strength}

\label{lesson:SA-C476-saktih}



\begin{quote}
Before the new one: say the Sanskrit for scratches, then the Sanskrit for a hiccup.
\end{quote}

\subsection*{You'll want to know: \sk{शक्तिः}}
\textbf{\sk{शक्तिः}} — \emph{śaktiḥ} — \textquotedblleft{}strength\textquotedblright{}.

\begin{quote}
\sk{व्यायामेन} \sk{शक्तिः} \sk{वर्धते।}
\end{quote}

\begin{quote}
\emph{vyāyāmena śaktiḥ vardhate.}
\end{quote}

\begin{quote}
Exercise builds strength.
\end{quote}

A feminine noun.

\begin{cousinweb}
\textbf{\sk{शक्तिः}} (\emph{śaktiḥ}) comes from the root √śak, \textquotedblleft{}to be able\textquotedblright{}, the root of \textbf{\sk{शक्नोति}}; MW gives \textquotedblleft{}power, ability, strength\textquotedblright{} (MW).
\end{cousinweb}

\begin{grammarlens}[title={What you've built}]
One more word for describing what you lived and felt.
\end{grammarlens}

\subsection*{Your turn}
\begin{itemize}
\raggedright
  \item \emph{Say it:} \emph{śaktiḥ}
  \item \emph{Say it:} \emph{śaktiḥ}, once more
  \item \emph{Say it:} \emph{vyāyāmena śaktiḥ vardhate.} — Exercise builds strength.
  \item \emph{Recall:} say the Sanskrit for scratches, then the Sanskrit for a hiccup, then say \emph{śaktiḥ} again
\end{itemize}

\begin{tcolorbox}[breakable,colback=teal!4,colframe=teal!35!black,title={Before you move on}]
What does \sk{शक्तिः} mean? (\textquotedblleft{}Strength\textquotedblright{}.) Say it once more.
\end{tcolorbox}
\section[\texorpdfstring{\sk{प्रांशुः} (prāṃśuḥ)}{प्रांशुः (prāṃśuḥ)}]{\sk{प्रांशुः} (prāṃśuḥ) — tall of stature}

\label{lesson:SA-C476-pramsuh}



\begin{quote}
Before the new one: say the Sanskrit for a hiccup, then the Sanskrit for strength.
\end{quote}

\subsection*{You'll want to know: \sk{प्रांशुः}}
\textbf{\sk{प्रांशुः}} — \emph{prāṃśuḥ} — \textquotedblleft{}tall of stature\textquotedblright{}.

\begin{quote}
\sk{मम} \sk{भ्राता} \sk{प्रांशुः} \sk{अस्ति।}
\end{quote}

\begin{quote}
\emph{mama bhrātā prāṃśuḥ asti.}
\end{quote}

\begin{quote}
My brother is tall.
\end{quote}

\begin{cousinweb}
\textbf{\sk{प्रांशुः}} (\emph{prāṃśuḥ}) MW reports that it is said to be built from \emph{pra} and \emph{aṃśu}, \textquotedblleft{}a ray\textquotedblright{}, so its origin is not settled; it gives \textquotedblleft{}high, tall, long\textquotedblright{} (MW).
\end{cousinweb}

\begin{grammarlens}[title={What you've built}]
One more word for describing what you lived and felt.
\end{grammarlens}

\subsection*{Your turn}
\begin{itemize}
\raggedright
  \item \emph{Say it:} \emph{prāṃśuḥ}
  \item \emph{Say it:} \emph{prāṃśuḥ}, once more
  \item \emph{Say it:} \emph{mama bhrātā prāṃśuḥ asti.} — My brother is tall.
  \item \emph{Recall:} say the Sanskrit for a hiccup, then the Sanskrit for strength, then say \emph{prāṃśuḥ} again
\end{itemize}

\begin{tcolorbox}[breakable,colback=teal!4,colframe=teal!35!black,title={Before you move on}]
What does \sk{प्रांशुः} mean? (\textquotedblleft{}Tall of stature\textquotedblright{}.) Say it once more.
\end{tcolorbox}
\section[\texorpdfstring{\sk{पुष्टः} (puṣṭaḥ)}{पुष्टः (puṣṭaḥ)}]{\sk{पुष्टः} (puṣṭaḥ) — well-nourished}

\label{lesson:SA-C476-pustah}



\begin{quote}
Before the new one: say the Sanskrit for strength, then the Sanskrit for tall of stature.
\end{quote}

\subsection*{You'll want to know: \sk{पुष्टः}}
\textbf{\sk{पुष्टः}} — \emph{puṣṭaḥ} — \textquotedblleft{}well-nourished\textquotedblright{}.

\begin{quote}
\sk{पुष्टः} \sk{वत्सः} \sk{क्षेत्रे} \sk{क्रीडति।}
\end{quote}

\begin{quote}
\emph{puṣṭaḥ vatsaḥ kṣetre krīḍati.}
\end{quote}

\begin{quote}
The well-fed calf plays in the field.
\end{quote}

\begin{cousinweb}
\textbf{\sk{पुष्टः}} (\emph{puṣṭaḥ}) is built on the root √puṣ, \textquotedblleft{}to thrive\textquotedblright{}; MW gives \textquotedblleft{}nourished, well-fed, thriving, strong\textquotedblright{} (MW).
\end{cousinweb}

\begin{grammarlens}[title={What you've built}]
One more word for describing what you lived and felt.
\end{grammarlens}

\subsection*{Your turn}
\begin{itemize}
\raggedright
  \item \emph{Say it:} \emph{puṣṭaḥ}
  \item \emph{Say it:} \emph{puṣṭaḥ}, once more
  \item \emph{Say it:} \emph{puṣṭaḥ vatsaḥ kṣetre krīḍati.} — The well-fed calf plays in the field.
  \item \emph{Recall:} say the Sanskrit for strength, then the Sanskrit for tall of stature, then say \emph{puṣṭaḥ} again
\end{itemize}

\begin{tcolorbox}[breakable,colback=teal!4,colframe=teal!35!black,title={Before you move on}]
What does \sk{पुष्टः} mean? (\textquotedblleft{}Well-nourished\textquotedblright{}.) Say it once more.
\end{tcolorbox}
\section[\texorpdfstring{\sk{आकृतिः} (ākṛtiḥ)}{आकृतिः (ākṛtiḥ)}]{\sk{आकृतिः} (ākṛtiḥ) — a figure; build}

\label{lesson:SA-C476-akrtih}



\begin{quote}
Before the new one: say the Sanskrit for tall of stature, then the Sanskrit for well-nourished.
\end{quote}

\subsection*{You'll want to know: \sk{आकृतिः}}
\textbf{\sk{आकृतिः}} — \emph{ākṛtiḥ} — \textquotedblleft{}a figure; build\textquotedblright{}.

\begin{quote}
\sk{तस्य} \sk{आकृतिः} \sk{पितुः} \sk{इव} \sk{अस्ति।}
\end{quote}

\begin{quote}
\emph{tasya ākṛtiḥ pituḥ iva asti.}
\end{quote}

\begin{quote}
He is built like his father.
\end{quote}

A feminine noun.

\begin{cousinweb}
\textbf{\sk{आकृतिः}} (\emph{ākṛtiḥ}) is built from \emph{ā} and the root √kṛ, \textquotedblleft{}to make\textquotedblright{}: how a thing is made; MW gives \textquotedblleft{}form, figure, shape, appearance\textquotedblright{} (MW).
\end{cousinweb}

\begin{grammarlens}[title={What you've built}]
One more word for describing what you lived and felt.
\end{grammarlens}

\subsection*{Your turn}
\begin{itemize}
\raggedright
  \item \emph{Say it:} \emph{ākṛtiḥ}
  \item \emph{Say it:} \emph{ākṛtiḥ}, once more
  \item \emph{Say it:} \emph{tasya ākṛtiḥ pituḥ iva asti.} — He is built like his father.
  \item \emph{Recall:} say the Sanskrit for tall of stature, then the Sanskrit for well-nourished, then say \emph{ākṛtiḥ} again
\end{itemize}

\begin{tcolorbox}[breakable,colback=teal!4,colframe=teal!35!black,title={Before you move on}]
What does \sk{आकृतिः} mean? (\textquotedblleft{}A figure; build\textquotedblright{}.) Say it once more.
\end{tcolorbox}
\section[\texorpdfstring{\sk{पोषयति} (poṣayati)}{पोषयति (poṣayati)}]{\sk{पोषयति} (poṣayati) — nourishes}

\label{lesson:SA-C476-posayati}



\begin{quote}
Before the new one: say the Sanskrit for well-nourished, then the Sanskrit for a figure.
\end{quote}

\subsection*{You'll want to know: \sk{पोषयति}}
\textbf{\sk{पोषयति}} — \emph{poṣayati} — \textquotedblleft{}nourishes\textquotedblright{}.

\begin{quote}
\sk{माता} \sk{शिशुं} \sk{क्षीरेण} \sk{पोषयति।}
\end{quote}

\begin{quote}
\emph{mātā śiśuṃ kṣīreṇa poṣayati.}
\end{quote}

\begin{quote}
The mother nourishes her baby with milk.
\end{quote}

Past (she …): \sk{सा} \sk{शिशुं} \sk{पोषितवती}, \emph{sā śiśuṃ poṣitavatī}.

\begin{cousinweb}
\textbf{\sk{पोषयति}} (\emph{poṣayati}) is built on the causative of the root √puṣ, \textquotedblleft{}to thrive\textquotedblright{}: to make thrive; MW gives \textquotedblleft{}to rear, nourish, feed\textquotedblright{} (MW).
\end{cousinweb}

\begin{grammarlens}[title={What you've built}]
One more word for describing what you lived and felt.
\end{grammarlens}

\subsection*{Your turn}
\begin{itemize}
\raggedright
  \item \emph{Say it:} \emph{poṣayati}
  \item \emph{Say it:} \emph{poṣayati}, once more
  \item \emph{Say it:} \emph{mātā śiśuṃ kṣīreṇa poṣayati.} — The mother nourishes her baby with milk.
  \item \emph{Recall:} say the Sanskrit for well-nourished, then the Sanskrit for a figure, then say \emph{poṣayati} again
\end{itemize}

\begin{tcolorbox}[breakable,colback=teal!4,colframe=teal!35!black,title={Before you move on}]
What does \sk{पोषयति} mean? (\textquotedblleft{}Nourishes\textquotedblright{}.) Say it once more.
\end{tcolorbox}
